% ==============================================================================
\section{Conclusions}
\label{sec:conclusions}
% ==============================================================================

We have introduced the aircraft hangar positioning problem (AHPP) and
provided a formal graph-based model centred on directed blocking arcs, a
compact mixed-integer linear programming formulation expressed at the
aircraft level (each aircraft as a single uninterrupted maintenance
block) that encodes entry and exit movements through a six-binary
indicator structure, and a suite of complementary heuristic methods
built on top of the same data model: a topology-aware GRASP that scores
positions by their blocking load, a fixed-assignment scheduler (FAS)
that re-optimises timing given a position assignment, and a safe
pipeline composition that returns the best of the components at no
extra cost.

The computational study on the 120-instance benchmark gives a clear
picture of when each method dominates.  Within the 60-second wall-clock
budget, the exact MILP only proves optimality on the small $R{=}5$
configuration; for every $R{=}10$ configuration the Gurobi gap remains
above $86\,\%$ --- $86\,\%$ even on the unconstrained \emph{None}
topology and above $93\,\%$ elsewhere --- which points to bound weakness
rather than blocking interactions as the dominant source of difficulty.
At $R{=}20$ the solver exhausts the budget with a gap close to
$100\,\%$, and at $R{=}30$ only a single seed returns a feasible
incumbent, so the heuristic pipeline becomes the primary source of
solutions for the largest configurations.  On
the $R{=}10$ configurations with non-trivial blocking the heuristic
pipeline cuts the best feasible MILP objective by $17$--$34\,\%$, and at
$R{=}20$ the gap widens by an order of magnitude (e.g.\ $\bar f = 681$
versus $2449$ on \emph{Triangle}, and $959$ versus $3332$ on
\emph{Full}).  The topology heuristic alone already matches the
safe-pipeline objective on every $R{=}10$ row, so the additional FAS
pass acts mostly as an insurance policy; its value becomes apparent at
the boundary of tractability, where FAS occasionally times out --- on
\emph{Triangle $R{=}30$}, for instance, FAS lags Topo by about $0.7\,\%$ and
the safe pipeline correctly falls back to the Topo solution.  The three
weight profiles act as effective dials between the objectives: the
default movement-priority profile returns \emph{zero} blocking
movements on every configuration covered by the benchmark at the cost
of a moderate makespan penalty; the delay-priority profile flattens
mean delay onto a narrow band on the $R{=}10$ rows while injecting tens
of movements; and the makespan-priority profile yields the lowest
$\bar m$ throughout, again at the cost of intermediate values of delay
and movements.

These results outline a natural agenda for future work.  A tighter MILP
formulation that exploits the capacity-one structure of the positions
and tightens the big-$M$ relaxation should close part of the bound gap
that currently bottlenecks the exact solver from $R{=}10$ onwards.  An
adaptive large neighbourhood search built on top of FAS, with destroy
and repair operators selected from observed per-instance-type
performance, is a promising avenue for the $R{\ge}20$ regime where FAS
already starts to feel the time limit.  A multi-shift extension that
allows aircraft arrivals to span several working days, and a validation
on production data from an airline maintenance partner, would complete
the bridge from the controlled benchmark presented here to operational
deployment.
